\documentclass[11pt]{article}
\usepackage[margin=1in]{geometry}
\usepackage{amsmath,amssymb,amsthm,graphicx,microtype}
\usepackage[colorlinks=true,linkcolor=blue,citecolor=blue,urlcolor=blue]{hyperref}
\newtheorem{theorem}{Theorem}
\newtheorem{remark}{Remark}
\newcommand{\Sone}{\mathbb S^1}
\newcommand{\PsiDO}{\Psi}
\title{A smoothing-unitary counterexample to a claimed\\spectral inequality for elliptic pseudodifferential operators}
\author{Automated mathematics run \texttt{fa\_banach\_001}, lane 2}
\date{11 August 2026}

\begin{document}
\maketitle

\noindent\textbf{Status.} Candidate counterexample, likely valid, awaiting
human review.  The claim inspected is Theorem~1.4 of Cardona
\cite{Cardona2209}, together with the principal-symbol interpretation of the
second open problem on page~18.

\section{Source problem and claimed theorem}

The source asks in Section~3 (physical page~18) for which manifolds the
restriction $\rho\geq 1-\delta$ can be removed, and conjectures that on a
manifold with symmetries, under a suitable positivity condition on the symbol,
the spectral inequality should hold throughout $0\leq\delta<\rho\leq1$.
The same paper's Theorem~1.4 claims, in the standard range, that every positive
elliptic $E\in\PsiDO^\nu_{\rho,\delta}(M)$ with positive symbol satisfies
\[
 \|v\|_{L^2(M)}\leq C_1e^{C_2\lambda}\|v\|_{L^2(\omega)}
 \quad
 \text{for }v\in\operatorname{span}\{\rho_j:\lambda_j\leq\lambda\}.
\]
In the proof, the relevant $E(x,\xi)$ is explicitly called the principal
symbol.  The following construction contradicts this theorem already for
$(\rho,\delta)=(1,0)$ and shows that principal-symbol positivity cannot be the
``suitable positivity condition'' sought by the open problem.

\begin{figure}[ht]
  \centering
  \includegraphics[width=\textwidth]{figures/open_problem_crop.png}
  \caption{The two open problems in Section~3 of arXiv:2209.10690,
  rendered from physical page~18.  The second bullet is the selected target.}
\end{figure}

\section{Proof intuition}

Ellipticity and the principal symbol only control an operator modulo smoothing
terms.  Smoothing terms can nevertheless rearrange individual smooth
eigenfunctions completely.  Start with the standard positive elliptic operator
$A=(I+L)^{\nu/2}$ on the circle, whose constant eigenfunction is everywhere
visible.  A rank-one Householder reflection $U=I+R$, with $R$ smoothing, can
send that constant to any prescribed smooth bump $f$ supported away from the
sensor.  The conjugate $E=UAU^*$ is still positive elliptic and has exactly the
same principal symbol as $A$, but $f$ is now an eigenfunction.  The proposed
observability estimate then reads $1\leq0$.

\section{Formal counterexample}

\begin{theorem}\label{thm:counterexample}
Let $\nu>0$.  There exist a nonempty open set $\omega\subset\Sone$ and a
positive, self-adjoint, invertible elliptic pseudodifferential operator
\[
 E\in\PsiDO^\nu_{1,0}(\Sone)
\]
whose principal symbol is the positive symbol $|\xi|^\nu$ for $\xi\ne0$, but
for which the spectral inequality of \cite[Theorem~1.4]{Cardona2209} fails.
Indeed, $E$ has a normalized $C^\infty$ eigenfunction $f$ satisfying
$f|_\omega=0$.
\end{theorem}

\begin{proof}
Give $\Sone$ normalized Haar measure and let
$L=-d^2/dx^2\geq0$.  Put
\[
 A=(I+L)^{\nu/2}.
\]
By the standard functional calculus, $A$ is positive, self-adjoint,
invertible, belongs to $\PsiDO^\nu_{1,0}(\Sone)$, and has principal symbol
$|\xi|^\nu$.  Its normalized constant function $h\equiv1$ satisfies $Ah=h$.

Choose a proper open arc $\omega$ whose complement has nonempty interior.
Choose a real $f\in C^\infty(\Sone)$ with
\[
 \operatorname{supp}f\subset\Sone\setminus\omega,
 \qquad \|f\|_2=1.
\]
Work first on the real Hilbert space.  Since $f\ne h$, define
\[
 w=\frac{h-f}{\|h-f\|_2},\qquad
 P_wu=\langle u,w\rangle w,\qquad U=I-2P_w.
\]
The operator $U$ is a self-adjoint unitary.  If
$c=\langle h,f\rangle\in\mathbb R$ and
$d=\|h-f\|_2$, then $d^2=2(1-c)$ and
$\langle h,w\rangle=d/2$, $\langle f,w\rangle=-d/2$.  Consequently
\[
 Uh=h-dw=f,\qquad Uf=f+dw=h.
\]
Extend $U$ complex linearly to complex $L^2(\Sone)$.

Now set
\[
 E=UAU.
\]
Unitary conjugation shows that $E$ is positive, self-adjoint, invertible, and
unitarily equivalent to $A$.  It remains to check the microlocal assumptions.
The rank-one projection $P_w$ has smooth Schwartz kernel
$w(x)\overline{w(y)}$, hence $P_w\in\PsiDO^{-\infty}$.  Expanding gives
\[
 E-A=-2P_wA-2AP_w+4P_wAP_w.
\]
Every term on the right has a smooth finite-rank kernel: for example, the
kernels of $AP_w$ and $P_wA$ are respectively
$(Aw)(x)\overline{w(y)}$ and
$w(x)\overline{(Aw)(y)}$.  Thus $E-A\in\PsiDO^{-\infty}$.
It follows that $E\in\PsiDO^\nu_{1,0}$ is elliptic and has the same positive
principal symbol as $A$.  Also $U=I+\PsiDO^{-\infty}$ preserves all Sobolev
spaces, so the unbounded-operator domain causes no change of class.

Finally,
\[
 Ef=UAUf=UAh=Uh=f.
\]
Thus $f$ is an eigenfunction of eigenvalue $1=1^\nu$ and belongs to the
spectral subspace at threshold $\lambda=1$.  The asserted inequality would
give
\[
 1=\|f\|_{L^2(\Sone)}
 \leq C_1e^{C_2}\|f\|_{L^2(\omega)}=0,
\]
a contradiction.
\end{proof}

\begin{remark}[Null controllability]
The same invisible eigenmode also obstructs the claimed null controllability:
the observation/control multiplier $1_\omega$ annihilates $f$, so the adjoint
observability inequality fails on the one-dimensional invariant subspace
spanned by $f$.
\end{remark}

\section{Verification and scope}

The construction is exact and uses no computation.  The decisive checks are:
(i) a smooth rank-one kernel is smoothing; (ii) smoothing perturbations preserve
the principal symbol and ellipticity; (iii) unitary conjugation preserves
positivity and spectrum; and (iv) the Householder identities swap $h$ and $f$.
All four are displayed explicitly in the proof.

This result does not solve the source's first, boundary/polyharmonic open
problem.  It also does not refute the deliberately unspecified existential
phrase ``under a suitable positivity condition.''  Rather, it proves sharply
that the hypotheses of source Theorem~1.4---positive operator and positive
principal symbol---are insufficient.  The later compact-Lie-group theorem of
Cardona, Delgado and Ruzhansky \cite{CDR2209} assumes nonnegativity of the full
global matrix-valued symbol.  Our construction is not claimed to satisfy that
stronger, quantization-specific condition.

\section{Attempt history and novelty check}

An initial route observed that Navier powers of the Dirichlet Laplacian inherit
the Laplacian spectral inequality, but that proper subcase does not reach the
clamped boundary conditions in the first open problem.  The smoothing-unitary
route upgrades the investigation to the complete counterexample above.

The bounded search on 11 August 2026 used the exact source title and arXiv id
with the terms \emph{correction}, \emph{erratum}, and \emph{counterexample},
and searched close variants involving spectral inequalities, elliptic
pseudodifferential operators, polyharmonic boundary conditions, and
Lebeau--Robbiano estimates.  No correction or matching counterexample was
found.  The close works arXiv:2209.12092 \cite{CDR2209} and arXiv:2309.02181
were inspected.  Novelty confidence is moderate because finite-rank smoothing
conjugation is standard folklore even though this exact application was not
located.

\section{Human review recommendation}

A microlocal-analysis reviewer should focus on the source's interpretation of
$E(x,\xi)>0$.  Its proof explicitly identifies this as the principal symbol;
under that reading Theorem~\ref{thm:counterexample} is complete.  If a stronger
complete-symbol condition was intended, it must be stated separately and the
scope of the source theorem changes.

\begin{thebibliography}{9}
\bibitem{Cardona2209}
D. Cardona,
\emph{Spectral inequalities for elliptic pseudo-differential operators on closed manifolds},
arXiv:2209.10690 (2022).

\bibitem{CDR2209}
D. Cardona, J. Delgado, and M. Ruzhansky,
\emph{Estimates for sums of eigenfunctions of elliptic pseudo-differential operators on compact Lie groups},
arXiv:2209.12092 (2022; revised 2022).

\bibitem{Hormander}
L. H\"ormander,
\emph{The Analysis of Linear Partial Differential Operators III},
Springer, 1985.
\end{thebibliography}

\end{document}

